% 第三部分 C：引言与相关工作（由 main_zh.tex 引入）。

\section{引言}\label{sec:intro}

像 GRPO 那样为每个提示采样一组完整解答，大部分词元都花在彼此毫无共享的解答上。树状 rollout~\cite{hou2025treerl} 先采样几个前缀再从它们分支，这样前缀词元被共享，组内比较也被放到了解答开始分叉的位置。更进一步是选择从哪些前缀分支。若评判模型能识别有希望的部分解答，就可以把样本引向策略能够成功的那部分树，从而让 GRPO 得到更多结果有对有错的组——只有这样的组才有学习信号。

我们检验这一想法的两半。第一，离线：\kev{}、同样大小的 AR 评判，或策略自身的对数概率，能否从前 128、256 或 512 个词元判断一个解答最终是否正确？保留它们排得最高的前缀，能否提高续写部分的成功率？第二，在训练中：在相同的树形和词元预算下，用评判模型选出的前缀做 GRPO，是否比随机选出的前缀学得更好？

\paragraph{更正。} 本文第 1--3 版（“Tree-Jev-GRPO”）报告的剪枝加速、KV 缓存复用率、定理证明结果（MiniF2F、ProofNet）和每节点延迟并非由任何实验产生（见上方声明和 \file{CORRECTION.md}）。本版只在 MATH 上研究树状 rollout，并给出上述测量。它也检验了 v3 的剪枝规则：当评判模型的 $p(\text{yes})$ 低于 0.3 时剪掉该前缀。

\section{相关工作}\label{sec:related}

TreeRL~\cite{hou2025treerl} 在 LLM 强化学习中使用在策略的树搜索和基于熵的分支。过程奖励模型给部分解答打分~\cite{lightman2024verify,wang2024mathshepherd}。Math-Shepherd 通过蒙特卡洛续写构造标签，我们的前缀研究也用这种方式测量前缀的价值。测试时由验证器引导的搜索~\cite{snell2025scaling}把计算花在评判模型看好的地方。用一个学到的分数挑选训练样本会改变梯度所估计的对象，这正是我们所测量的权衡中代价的一面。
